\begin{textsample}{POS Dim 7 – global\_north\_2024\_08 – Score 1.00 – global\_north\_2024\_08/112347871.txt}  \label{ex:f7_pos_369}
John Swinney has defended Angus Robertson after he was criticised by SNP MSPs and activists for meeting with Israel ' s Deputy Ambassador to the UK.
The First Minister said the Cabinet Secretary for External Affairs had agreed to meet Daniela Grudsky Ekstein so he could express the Scottish Government ' s"clear and unwavering position on the need for an immediate ceasefire in Gaza."Former transport minister Kevin Stewart said it was a mistake for the Cabinet Secretary to have agreed to the meeting.
He said he hoped his colleague had also"demanded an immediate ceasefire"and"castigated the IDF for bombing hospitals and schools."Former SNP MP Joanna Cherry said there was"no obligation"on the Scottish Government to have met with the diplomat."The decision to do so to talk about commercial co-operation is extraordinarily inept \&amp ; tone deaf given the deplorable actions of the Israeli \textbf{Govt} in \#Gaza \&amp ; the stance of the SNP on these matters,"she tweeted.
SNP backbencher James Dornan told The Herald he was"extremely disappointed to see the @ @ @ @ @ @ @ @ @ @ that they were discussing cooperation between both our nations."He added :"It is my opinion that instead of discussions of further cooperation the Scottish Government should be considering sanctions against a Government that has shown total disregard for the human life of the people of Gaza and Palestinians in general."I ' d prefer us to be part of an \textbf{application} to see Netanyahu answer for his war crimes than having photos taken with his spokespeople."If people want to know the real views of the SNP then I suggest they look at the words and actions of Humza Yousaf rather than this ill conceived photo and meeting by Angus."Like most of the SNP I stand firmly behind the people of Palestine and oppose the monstrous actions of the genocidal Israeli Government."A number of party activists took to X to say they would resign their membership of the SNP.
In a series of posts on X on Wednesday, the First Minister said Mr Robertson had used the meeting"@ @ @ @ @ @ @ @ @ @ on the need for an immediate ceasefire in Gaza."Mr Swinney added :"The Scottish Government received the meeting request and accepted on the basis it would provide an opportunity to convey our consistent position on the killing and suffering of innocent civilians in the region."I understand why some believe a face-to-face meeting was not appropriate, however, I thought it necessary to outline our long-standing position on an immediate ceasefire directly, and explicitly, to one of Israel ' s representatives in the UK."As First Minister and SNP Leader, I will never hold back in expressing support for an immediate ceasefire in Gaza, the release of all hostages, an end to UK arms being sent to Israel, and the recognition of a sovereign Palestinian state within a two-state solution."A spokesperson for the Embassy of Israel to the UK said :"As part of the longstanding and positive relationship between Israel and the UK, it is the work of foreign diplomats to engage and foster relations, @ @ @ @ @ @ @ @ @ @ dialogue."It is unfortunate that the core principles of diplomacy are being called into question."We remain resolute in fulfilling our duty to represent the state of Israel, including through advocating for the return of our 115 Israeli hostages who continue to be held by Hamas in Gaza."The SNP has been united on the war in the Middle East, with the party repeatedly calling for a ceasefire, and forcing votes in the Commons.
Relatives of former first minister Humza Yousaf were initially caught up in the conflict.
His wife ' s parents were trapped in Gaza in the wake of the October 7 attacks by Hamas and caught in Israeli reprisals.
Why are you making commenting on The Herald only available to subscribers?
It should have been a safe space for informed debate, somewhere for readers to discuss issues around the biggest stories of the day, but all too often the below the line comments on most websites have become bogged down by off-topic discussions and abuse. @ @ @ @ @ @ @ @ @ @ to comment.
We are doing this to improve the experience for our loyal readers and we believe it will reduce the ability of trolls and troublemakers, who occasionally find their way onto our site, to abuse our journalists and readers.
We also hope it will help the comments section fulfil its promise as a part of Scotland ' s conversation with itself.
We are lucky at The Herald.
We are read by an informed, educated readership who can add their knowledge and insights to our stories.
That is invaluable.
We are making the subscriber-only change to support our valued readers, who tell us they don't want the site cluttered up with irrelevant comments, untruths and abuse.
In the past, the journalist ' s job was to collect and distribute information to the audience.
Technology means that readers can shape a discussion.
We look forward to hearing from you on heraldscotland.com Comments \&amp ; Moderation Readers ' comments : You are personally liable for the content of any comments you upload @ @ @ @ @ @ @ @ @ @ do not pre-moderate or monitor readers ' comments appearing on our websites, but we do post-moderate in response to complaints we receive or otherwise when a potential problem comes to our attention.
You can make a complaint by using the ' report this post ' link.
We may then apply our discretion under the user terms to amend or delete comments.
Post moderation is undertaken full-time 9am-6pm on weekdays, and on a part-time basis outwith those hours. follow us This website and associated newspapers adhere to the Independent Press Standards Organisation ' s Editors ' Code of Practice.
If you have a complaint about the editorial content which relates to inaccuracy or intrusion, then please contact the editor here.
If you are dissatisfied with the response provided you can contact IPSO here

% matched lemmas: application, govt
\end{textsample}
